\chapter{Introduction}

\section{Motivation}

In the past five decades, the world has undergone a demographic transformation of unprecedented scale and pace. The global urban population has grown from fewer than one billion people in 1960 to more than four billion today, a trajectory projected to intensify throughout the remainder of the twenty-first century as populations across Asia, Africa, and Latin America continue migrating toward metropolitan centers in pursuit of economic opportunity, social mobility, and access to public services\footnote{\label{fn:UNH} UN-Habitat (2020). \textit{World Cities Report 2020: The Value of Sustainable Urbanization.} \url{https://unhabitat.org/world-cities-report-2020-the-value-of-sustainable-urbanization}}. This concentration of human activity within urban boundaries has made cities simultaneously the principal engines of global economic productivity and among the principal sites of environmental pressure: they account for approximately 70 percent of total energy consumption, generate a commensurate share of greenhouse gas emissions, and place sustained demand on water, transportation, and sanitation infrastructure. The capacity to accurately monitor, quantify, and interpret the spatial dynamics of urban growth has therefore emerged as a foundational requirement for sustainable development planning, disaster risk reduction, and evidence-based governance at all scales of public administration.

Lima Metropolitan Area clearly exemplifies the tensions characterizing urban expansion in the Global South. Home to more than ten million inhabitants concentrated in one of the driest coastal desert environments on Earth and situated along one of the world's most seismically active subduction zones, Lima's built fabric has grown through a dual process of horizontal extension and vertical intensification\footnote{Instituto Nacional de Estadística e Informática (INEI), Censo Nacional de Población 2017. \url{https://censo2017.inei.gob.pe/resultados-definitivos-de-los-censos-nacionales-2017/}}. Horizontal expansion refers to the occupation of previously undeveloped land at the urban periphery (Figure~\ref{fig:horizontal}): hillside settlements constructed with no formal engineering oversight, informal subdivisions extending across desert flatlands, and the gradual absorption of adjacent agricultural and ecological zones into the consolidated metropolitan area \footnote{Water-Sensitive Urban Plan for Lima Metropolitan Area (Peru) Based on Changes in the Urban Landscape from 1990 to 2021. doi: \url{https://doi.org/10.3390/land11122261}}. Vertical densification, by contrast, refers to the progressive addition of stories to the existing building stock (Figure~\ref{fig:vertical}), a process that has been documented in Lima's established residential neighborhoods, where rising land values and in-situ population pressure compel owners to maximize floor area within existing footprints \footnote{Confederación Nacional de Instituciones Empresariales Privadas (2013, December 13). Residential buildings in Lima gained two floors in height since 2011. CONFIEP. \url{https://www.confiep.org.pe/noticias/edificios-de-viviendas-en-lima-ganaron-dos-pisos-en-altura-desde-el-2011/}}. Both modalities impose distinct demands on infrastructure networks, seismic load calculations, and public utility capacity, yet they differ fundamentally in their spectral and geometric signatures in satellite imagery (Figure~\ref{fig:comparison_growth}), making their automated discrimination a non-trivial remote sensing challenge. National census data, updated only once per decade, are insufficient to capture the pace of contemporary urban growth, while municipal cartographic records often lag behind observed physical development by several years, leaving planning agencies without timely, high-resolution spatial information required for effective urban planning.

\begin{figure}[h!]
    \centering
    \begin{subfigure}[b]{0.45\textwidth}
        \centering
        \includegraphics[height=6cm]{graficos/horizontal.png}
        \caption{Urban horizontal expansion: progressive occupation of peripheral land by new residential and commercial settlements. \textbf{Source:} \textit{Water-Sensitive Urban Plan for Lima Metropolitan Area (Peru) Based on Changes in the Urban Landscape from 1990 to 2021}. DOI: \url{https://doi.org/10.3390/land11122261}.}
        \label{fig:horizontal}
    \end{subfigure}
    \hfill
    \begin{subfigure}[b]{0.45\textwidth}
        \centering
        \includegraphics[height=6cm]{graficos/vertical.png}
        \caption{Urban vertical densification within consolidated urban areas. \textbf{Source:} Confederación Nacional de Instituciones Empresariales Privadas (CONFIEP, 2013), \textit{Residential buildings in Lima gained two floors in height since 2011} (\href{https://www.confiep.org.pe/noticias/edificios-de-viviendas-en-lima-ganaron-dos-pisos-en-altura-desde-el-2011/}{https://confiep.org.pe/noticias/edificios}).}
        \label{fig:vertical}
    \end{subfigure}

    \caption{The two primary modalities of urban growth observable in the Lima Metropolitan Area. An automated change detection system must discriminate between these structurally distinct patterns to provide decision-relevant intelligence for urban planning and infrastructure management.}
    \label{fig:comparison_growth}
\end{figure}

\ac{VHR} satellite imagery offers the most promising technological pathway for addressing this monitoring gap at scale. Contemporary Earth observation missions, including the Peruvian satellite PeruSAT-1, are capable of acquiring panchromatic imagery at 0.70 meters per pixel and multispectral imagery at 2.80 meters per pixel, thereby resolving individual building footprints, rooftop configurations, and construction material characteristics across large geographic extents within a single orbital pass. By comparing co-registered image pairs acquired at different temporal epochs, automated \ac{BCD} algorithms can identify which structures have appeared, disappeared, expanded, or undergone material modification between acquisition dates, a capability of direct operational value for urban planning, informal settlement monitoring, post-disaster damage assessment, and infrastructure auditing \cite{Cheng2021DeepLab, Sirko2023TeachStudentSentinel2}. The spatial detail required for this task is only accessible to \ac{VHR} sensors: as shown in Figure~\ref{fig:resolution_comparison}, individual building footprints that PeruSAT-1 resolves at 0.70~m per pixel are entirely lost at the medium resolution of sensors such as Sentinel-2. The central methodological challenge, however, is that the volumes of imagery generated by modern satellite constellations far exceed the annotation and manual analysis capacity of expert practitioners, and the heterogeneous radiometric conditions, geometric misregistration, and seasonal spectral variation inherent in multitemporal archives severely complicate straightforward image-differencing approaches. In recent years, deep learning has therefore emerged as the predominant paradigm for scalable, robust \ac{BCD} at city or regional scale \cite{Sirko2023TeachStudentSentinel2, Cheng2021DeepLab}.

\begin{figure}[htbp]
    \centering
    \begin{subfigure}[b]{0.48\linewidth}
        \centering
        \includegraphics[width=\linewidth, height=6.5cm]{graficos/VHR-sentinel2.png}
        \caption{Comparison of VHR imagery against medium-resolution Sentinel-2 (10~m/pixel): individual buildings cannot be resolved at medium resolution \cite{Sirko2023TeachStudentSentinel2}.}
    \end{subfigure}
    \hfill
    \begin{subfigure}[b]{0.48\linewidth}
        \centering
        \includegraphics[width=\linewidth, height=6.4cm]{graficos/PeruSAT1.png}
        \caption{PeruSAT-1 panchromatic imagery at 0.70~m/pixel. \textbf{Source:} Infoespacial (2016), \textit{``El PeruSAT-1 envía a la Tierra sus primeras imágenes''} [\href{https://www.infoespacial.com/texto-diario/mostrar/3568529/039perusat-1-039-envia-tierra-primeras-imagenes}{Info Espacial - PeruSat}].}
    \end{subfigure}
    \caption{Spatial resolution comparison between satellite imaging systems. VHR missions such as PeruSAT-1 are essential for building-level change detection, whereas medium-resolution sensors provide broad coverage without the spatial detail required to resolve structural modifications in the built environment.}
    \label{fig:resolution_comparison}
\end{figure}

The past decade has witnessed a marked paradigm shift in remote sensing change detection, driven by the representational capacity of deep learning architectures. Following the success of fully convolutional networks and encoder-decoder designs adapted for semantic segmentation, the introduction of attention mechanisms and \ac{ViT} backbones \cite{Vaswani2017Attention, Dosovitskiy2021ViT} opened a new frontier for modeling non-local spatial and temporal dependencies across entire image regions. Within the \ac{BCD} domain, architectures such as the Binary change detection with Image Transformers (BIT) \cite{Chen2022BIT}, SNUNet-CD \cite{Fang2021SNUNetCD}, and ChangeFormer \cite{Bandara2022ChangeFormer} demonstrated that globally context-aware, hierarchical Siamese encoders equipped with multi-head self-attention substantially outperform locally-receptive convolutional predecessors on challenging urban datasets, with ChangeFormer \cite{Bandara2022ChangeFormer} achieving state-of-the-art accuracy on the LEVIR-CD benchmark \cite{Zheng2020LEVIR}. These advances confirmed a generalized principle: accurate discrimination of building-scale structural changes requires simultaneous sensitivity to fine-grained local texture differences and broad spatial context \cite{Zheng2022HFANet, Teng2023SwinCD}, capable of disambiguating true construction events from radiometric artifacts, illumination variations, and seasonal phenological effects.

The representational advantages of Transformer-based \ac{BCD}, however, carry a computational cost that constitutes a critical accessibility barrier. Standard multi-head self-attention scales quadratically in both time and memory with respect to the number of input tokens a regime that becomes acutely limiting in \ac{VHR} remote sensing, where even moderate patch sizes generate hundreds of spatial embeddings, and bi-temporal pipelines further amplify this burden by processing pre-event and post-event images through shared encoder streams. Researchers operating outside high-performance computing environments a condition representative of most academic and public-sector institutions in Latin America and the Global South find that this memory footprint constrains feasible batch sizes, effective training resolutions, and iteration speed to a degree that impedes systematic model development and comparative experimentation. \ac{RALA} \cite{Fan2025RALA} constitutes a principled architectural response to this constraint: by reformulating the key-value aggregation step as a rank-augmented weighted buffer, it reduces attention-related peak \ac{GPU} memory consumption by approximately 22 percent and increases the maximum feasible training batch size by up to 50 percent at standard processing resolutions  efficiency gains that broaden access to Transformer-based \ac{BCD} for research groups operating under constrained hardware budgets. Integrating \ac{RALA} into the encoder of ChangeFormer therefore offers the opportunity to extend state-of-the-art change detection capacity to resource-limited environments, including those engaged in \ac{VHR} urban monitoring of Peruvian cities with PeruSAT-1 imagery.


\section{Problem}

The central architectural bottleneck addressed in this thesis arises from the computational properties of standard \ac{MHSA} in Transformer-based change detection networks. In \ac{MHSA}, each token in the input sequence attends to all other tokens by computing pairwise similarity scores, forming a dense interaction matrix whose size grows as the square of the sequence length \cite{Vaswani2017Attention}. In natural image classification tasks, where input sequences are relatively short and a single global representation suffices, this quadratic regime remains tractable. In bi-temporal \ac{VHR} building change detection, two compounding factors transform this scaling behavior into a decisive practical constraint: the spatial resolution of \ac{VHR} inputs generates long token sequences at each encoder level, and the inherently bi-temporal nature of the task effectively doubles the token count within each forward pass by processing pre-event and post-event image patches together. The result is a memory consumption profile that grows prohibitively as either input resolution or batch size increases a property that is particularly severe in hierarchical Siamese architectures such as ChangeFormer \cite{Bandara2022ChangeFormer}, where \ac{MHSA} operates across multiple scales simultaneously and token counts at finer encoder stages remain substantial.

The practical ramifications of this quadratic regime manifest directly in the experimental constraints facing researchers without access to high-memory \ac{GPU} clusters. Under a 40~GB \ac{VRAM} budget representative of modern academic computing environments a standard ChangeFormer model operating at $256 \times 256$ patch resolution supports a maximum batch size of 32 images per training iteration; scaling the patch resolution to $512 \times 512$ reduces this to 8 images, and operating at the native $1024 \times 1024$ resolution of the LEVIR-CD dataset reduces the feasible batch size to 2 a regime in which gradient estimation becomes noisy, convergence becomes unreliable, and the statistical diversity of training samples per update is severely diminished. These constraints are not merely inconveniences of scale: small effective batch sizes impair the representativeness of gradient updates, slow the learning trajectory on heterogeneous datasets, and make it impractical to evaluate the effect of input resolution on detection accuracy a parameter of particular relevance for PeruSAT-1 applications, where the 0.70~m per pixel panchromatic resolution demands high-resolution processing to resolve fine-grained structural changes such as rooftop modifications or single-storey additions. The combinatorial effect is a systematic reduction in research throughput and model exploration capacity for institutions operating without large-scale \ac{GPU} infrastructure.

An apparently natural remedy is to substitute the quadratic softmax attention kernel with a linear feature map, reducing the attention complexity from quadratic to linear in the sequence length. This family of efficient attention mechanisms has attracted substantial theoretical and empirical interest in the broader machine learning community. However, a critical failure mode has been formally identified for linear attention in dense prediction settings: the intermediate key-value aggregation buffer of linear attention tends toward rank collapse, a degeneration in which the aggregated sequence representation occupies a low-dimensional subspace that discards the spatial diversity essentialfor pixel-accurate segmentation and change localization. Earlier efficient attention methods demonstrated measurable accuracy degradation on dense vision benchmarks precisely for this reason, undermining their practical utility for change detection applications where precise boundary delineation is required. \ac{RALA} \cite{Fan2025RALA} directly addresses rank collapse through a learned, token-wise reweighting of key-value contributions augmented by a channel-wise output modulation step that restores feature diversity without reintroducing quadratic cost. The central research hypothesis motivating this thesis is therefore the following: when the rank-augmented linear attention mechanism of \ac{RALA} is systematically substituted for the standard \ac{MHSA} modules within the hierarchical encoder of ChangeFormer, can the resulting model maintain segmentation accuracy comparable to the quadratic baseline across diverse building change detection benchmarks spanning multiple acquisition geometries and scene complexities while simultaneously delivering meaningful reductions in \ac{GPU} memory consumption and enabling larger, more stable training configurations?


\section{Goals}

\noindent\textbf{Main Objective.} The primary objective of this thesis is to design, implement, and rigorously validate \textbf{RALA-ChangeFormer}, a computationally efficient hierarchical Siamese Transformer architecture for multi-temporal building change detection from \ac{VHR} satellite imagery. The proposed architecture pursues the replacement of standard quadratic self-attention within the encoder stages of ChangeFormer with rank-augmented linear attention \cite{Fan2025RALA}, with the dual ambition of reducing attention-related computational cost from quadratic to near-linear complexity while preserving the segmentation accuracy of the baseline model across diverse urban environments and acquisition conditions. This primary objective is motivated not only by the theoretical interest of studying efficiency-accuracy trade-offs in bi-temporal Transformer architectures, but also by the applied imperative of extending state-of-the-art change detection capabilities to research and operational contexts where access to large-memory \ac{GPU} infrastructure cannot be assumed including Peruvian institutions, such as the Urban Data Analysis subgroup within GINIA at UTEC, engaged in monitoring PeruSAT-1 imagery of Lima Metropolitan Area and other rapidly growing South
American cities.

\noindent\textbf{Secondary Objectives.} The realization of this primary objective is pursued through a coherent sequence of specific operational goals. The first concerns the rigorous architectural design and software implementation of the RALA-ChangeFormer encoder, requiring the adaptation of the rank-augmentation mechanism to the multi-scale, hierarchical token structure of ChangeFormer \cite{Bandara2022ChangeFormer} and the verification of tensor-dimensional compatibility, numerical stability, and accurate computational profiling across all encoder stages. The second goal is a comprehensive efficiency analysis that quantifies, under controlled experimental conditions, the reductions in \ac{FLOPs}, peak \ac{GPU} memory consumption, maximum feasible batch size, and per-epoch training throughput achieved by \ac{RALA} relative to the quadratic baseline providing empirical grounding for the argument that efficient attention mechanisms are a meaningful tool for broadening access to advanced \ac{BCD} research. The third, and most extensive, goal is the systematic cross-domain evaluation of RALA-ChangeFormer across six geographically and methodologically diverse benchmarks: the LEVIR-CD and DSIFN-CD datasets that constitute the foundational experimental suite, together with four extended datasets WHU-CD, LEVIR-CD+, SYSU-CD, and S2Looking  spanning a range of acquisition geometries from nadir-view centimetre-resolution aerial imagery to oblique side-looking satellite acquisitions, annotation quality levels from precisely delineated building footprints to multi-class surface change labels, and scene heterogeneity conditions from homogeneous planned urban developments to complex informal settlement environments. This extended experimental program is designed not merely to assess whether RALA-ChangeFormer achieves performance parity with the baseline under standard benchmark conditions, but to identify the generalization boundaries of rank-augmented linear attention: the acquisition characteristics that preserve the efficiency advantage without accuracy sacrifice, and the failure modes that motivate further architectural development.

\noindent\textbf{Outline.} The remainder of this document is organized to build progressively from conceptual foundations toward full experimental validation. Chapter~\ref{chap:background} establishes the theoretical prerequisites required to understand the proposed method, covering the principles of satellite imagery and spatial resolution, the formal definition of the bi-temporal building change detection problem, the mathematical structure of multi-head self-attention and its quadratic complexity, the rank-augmented linear attention formulation underpinning \ac{RALA}, and the evaluation metrics employed throughout this work. Chapter~\ref{chap:related_work} situates the contribution within the existing literature by reviewing supervised, semi-supervised, and self-supervised learning strategies for \ac{BCD}, the progression from convolutional Siamese networks to hierarchical Transformer architectures, and the landscape of efficient attention mechanisms and their known limitations in dense prediction settings. Chapter~\ref{chap:method} details the complete RALA-ChangeFormer pipeline, including the preprocessing stages applied to bi-temporal image pairs, the precise integration of rank-augmented attention into each encoder block, the hybrid loss function, and the post-processing localization procedure that converts model outputs into interpretable binary change masks. Chapter~\ref{chap:experiments} presents the full experimental evaluation, encompassing quantitative accuracy and efficiency comparisons across all six benchmarks, qualitative analysis of predicted change masks and attention distributions, training dynamics under each acquisition condition, and a substantive discussion of the performance gradient observed across datasets of varying geometric and spectral complexity. Finally, Chapter~\ref{chap:conclusiones} synthesizes the principal findings, enumerates the current limitations of the proposed approach, and outlines future research directions including geometric alignment for oblique-view imagery, dynamic rank adaptation strategies, and the construction of dedicated building change detection benchmarks for the heterogeneous urban environments of the Global South.


\section{Contributions}

The main contributions of this thesis, first introduced in our conference publication \cite{acuna2026efficient}, can be summarized as follows:

\noindent\textbf{A lightweight change detection model.} We propose RALA-ChangeFormer, a hierarchical Siamese Transformer that replaces the quadratic multi-head self-attention of ChangeFormer with rank-augmented linear attention, reducing attention-related peak \ac{GPU} memory consumption and enabling larger, more stable training configurations without sacrificing segmentation accuracy.

\noindent\textbf{A complete and robust benchmarking.} We conduct a systematic cross-domain evaluation of the proposed model across six geographically and methodologically diverse datasets (LEVIR-CD, DSIFN-CD, WHU-CD, LEVIR-CD+, SYSU-CD, and S2Looking), jointly reporting accuracy and efficiency metrics (FLOPs, peak memory, maximum feasible batch size, and training throughput) to characterize the generalization boundaries of rank-augmented linear attention.

\noindent\textbf{A peer-reviewed publication.} The core of this work has been accepted and published at the 21st International Conference on Computer Vision Theory and Applications (VISAPP 2026) \cite{acuna2026efficient}, validating the relevance and soundness of the proposed approach.

\noindent\textbf{An ongoing journal extension.} Building on the accepted conference paper, we are currently developing an extended version for a journal submission, broadening the experimental evaluation and deepening the analysis of the efficiency-accuracy trade-offs reported throughout this thesis.